\chapter{Seguimiento del proyecto y decisiones relevantes}
\label{anexo:seguimiento_decisiones}

Este anexo recoge las decisiones técnicas más relevantes adoptadas durante el desarrollo de TierraViva. Su finalidad es documentar cómo se ha ido acotando el proyecto, qué alternativas se han descartado y qué criterios han guiado la arquitectura final del sistema.

El anexo complementa la planificación temporal recogida en el Anexo \ref{anexo:cronograma}, los recursos humanos y experimentales descritos en el Anexo \ref{anexo:plan_recursos} y el análisis de riesgos, oportunidades y limitaciones generales desarrollado en el Anexo \ref{anexo:gestion_riesgos}. Su contenido se centra en las decisiones de ingeniería que condicionan el diseño del prototipo y explican la evolución técnica de la arquitectura.

\section{Criterio de registro}

Las decisiones se documentan atendiendo a cuatro aspectos:

\begin{itemize}
    \item \textbf{situación de partida}: problema o duda técnica que existía antes de tomar la decisión;
    \item \textbf{decisión adoptada}: elección realizada en el proyecto;
    \item \textbf{justificación}: motivo técnico, experimental o de alcance;
    \item \textbf{impacto}: consecuencia directa sobre la arquitectura, la implementación o las pruebas.
\end{itemize}

Este criterio permite mostrar que la solución final es el resultado de un proceso de selección, descarte y ajuste progresivo, en el que cada tecnología cumple una función concreta dentro de la arquitectura.

\section{Decisiones técnicas relevantes}

La Tabla \ref{tab:registro_decisiones_tecnicas} resume las decisiones que más han influido en la arquitectura final de TierraViva.

\begin{table}[H]
\centering
\renewcommand{\arraystretch}{1.22}
\scriptsize
\begin{tabular}{|p{0.09\textwidth}|p{0.24\textwidth}|p{0.29\textwidth}|p{0.28\textwidth}|}
\hline
\rowcolor{URJCRed}
\TVHead{ID} & \TVHead{Decisión} & \TVHead{Justificación} & \TVHead{Impacto en el proyecto} \\
\hline
D01 & Orientar TierraViva a riego inteligente basado en IoT. & Permitía combinar sensorización, electrónica, comunicaciones, software y aplicación ambiental real dentro de un alcance defendible para un TFG. & El proyecto queda centrado en un prototipo funcional de monitorización y actuación sobre riego. \\
\hline
D02 & Utilizar Arduino UNO como controlador de estación. & Es una plataforma suficiente para leer sensores, controlar el relé y comunicarse con el módem mediante serie, manteniendo baja la complejidad de la estación. & La estación se mantiene simple, reproducible y adecuada para pruebas por bloques. \\
\hline
D03 & Usar Raspberry Pi como nodo local de supervisión. & La Raspberry Pi aporta capacidad para ejecutar ingesta, SQLite, API y \textit{dashboard} desde un nodo común de supervisión. & Se separa la estación de campo del sistema de supervisión y visualización. \\
\hline
D04 & Emplear conectividad LTE mediante A7670E-LASA. & El sistema está pensado para parcelas o entornos donde resulta conveniente priorizar cobertura móvil frente a una red WiFi local. & La estación puede publicar telemetría usando infraestructura móvil. \\
\hline
D05 & Utilizar ThingSpeak como \textit{broker} de telemetría. & Permite recibir datos IoT publicados por las estaciones y consultarlos posteriormente desde Raspberry Pi sin desarrollar una plataforma \textit{cloud} propia. & Reduce tiempo de desarrollo y centra el esfuerzo en integración, validación y visualización. \\
\hline
D06 & Mantener la decisión de riego en Arduino. & La actuación sobre agua debe poder ejecutarse de forma segura aunque falle temporalmente la comunicación o la supervisión. & La estación conserva autonomía local para decidir la actuación y mantener la bomba apagada ante condiciones no válidas. \\
\hline
D07 & Mantener el \textit{dashboard} como herramienta de supervisión. & El control descendente requiere seguridad, sincronización, gestión de comandos, latencia controlada y prevención de actuaciones no deseadas. & El \textit{dashboard} presenta lecturas, estados y eventos publicados por las estaciones, mientras que la actuación física permanece en el \textit{firmware}. \\
\hline
D08 & Usar SQLite como base de datos local. & Es una solución ligera, basada en fichero y suficiente para el volumen de datos del prototipo. & Permite conservar histórico local en la Raspberry Pi con una arquitectura sencilla y reproducible. \\
\hline
D09 & Implementar la API con FastAPI. & Facilita exponer datos estructurados mediante HTTP y conectar el \textit{backend} con el \textit{dashboard}. & Se separan almacenamiento, lógica de consulta e interfaz web. \\
\hline
D10 & Crear un \textit{dashboard} web local. & Permite consultar el sistema desde un navegador mediante una interfaz servida por la Raspberry Pi. & Mejora la usabilidad y permite visualizar lecturas, estados, gráficas e histórico. \\
\hline
D11 & Validar el relé antes de conectar la bomba. & La actuación física introduce riesgos eléctricos e hidráulicos que conviene aislar durante las primeras pruebas. & Se establece una secuencia de validación más segura: sensores, comunicación, relé sin carga y bomba controlada. \\
\hline
D12 & Documentar limitaciones y decisiones descartadas. & El proyecto combina hardware, comunicaciones y software, por lo que es importante diferenciar lo implementado, lo validado y lo planteado como trabajo futuro. & La memoria gana trazabilidad y delimita con claridad el alcance real de la versión desarrollada. \\
\hline
D13 & Incorporar un condensador electrolítico de \(2200~\mu\mathrm{F}\) y un diodo en la rama de la bomba. & La bomba de agua actúa como carga de motor y puede introducir perturbaciones durante el arranque, parada o conmutación mediante relé. & Se mejora la estabilidad eléctrica del bloque de actuación y se documenta una medida de mitigación aplicada durante el montaje del prototipo. \\
\hline
\end{tabular}
\caption{Registro de decisiones técnicas relevantes.}
\label{tab:registro_decisiones_tecnicas}
\end{table}

\section{Alternativas consideradas y descartadas}

Durante el desarrollo se valoraron alternativas que ayudaron a definir la arquitectura final. La Tabla \ref{tab:alternativas_descartadas} resume las principales opciones analizadas, el motivo de descarte y la decisión adoptada.

\begin{table}[H]
\centering
\renewcommand{\arraystretch}{1.22}
\scriptsize
\begin{tabular}{|p{0.24\textwidth}|p{0.32\textwidth}|p{0.32\textwidth}|}
\hline
\rowcolor{URJCRed}
\TVHead{Alternativa} & \TVHead{Motivo de descarte} & \TVHead{Decisión final} \\
\hline
Uso exclusivo de WiFi & Dependía de disponer de una red local cercana a la parcela o huerto. & Se adopta conectividad LTE para las estaciones remotas. \\
\hline
LoRa/LoRaWAN como comunicación principal & Requería infraestructura adicional, \textit{gateway} o cobertura LoRaWAN adecuada, aumentando la complejidad del prototipo. & Se prioriza LTE por viabilidad de integración y disponibilidad de red móvil. \\
\hline
Raspberry Pi como controlador directo del riego & Obliga a mantener conexión fiable con la estación y complica la actuación segura ante fallos. & La decisión de riego queda en Arduino. \\
\hline
Control remoto desde \textit{dashboard} & Añade gestión de órdenes, autenticación, seguridad, sincronización y prevención de activaciones no deseadas. & El \textit{dashboard} se mantiene como herramienta de monitorización. \\
\hline
Base de datos externa o servidor remoto & Introduce dependencias de red, configuración y mantenimiento innecesarias para el prototipo. & Se emplea SQLite local en Raspberry Pi. \\
\hline
Aplicación móvil propia & Aumenta el alcance del desarrollo y desplaza esfuerzo desde la integración IoT hacia una capa adicional de interfaz. & Se implementa un \textit{dashboard} web accesible desde navegador. \\
\hline
Modelo predictivo o inteligencia artificial & Requiere histórico suficiente, entrenamiento y validación experimental prolongada. & Se emplea lógica basada en umbrales y se deja la predicción como línea futura. \\
\hline
PCB propia & Mejoraría robustez, pero exige diseño, fabricación y validación adicional fuera del alcance actual. & Se mantiene montaje de prototipo documentado mediante esquemas y tablas de conexión. \\
\hline
Producto estanco definitivo & Requiere diseño mecánico, conectores adecuados y pruebas de intemperie prolongadas. & Se plantea protección básica para pruebas y mejora futura de encapsulado. \\
\hline
\end{tabular}
\caption{Alternativas consideradas y descartadas durante el desarrollo.}
\label{tab:alternativas_descartadas}
\end{table}

\section{Cambios de enfoque}

Algunas decisiones modificaron el papel de varios bloques dentro del sistema y condicionaron la arquitectura final. La Tabla \ref{tab:cambios_enfoque} recoge los cambios de enfoque más relevantes.

\begin{table}[H]
\centering
\renewcommand{\arraystretch}{1.22}
\small
\begin{tabular}{|p{0.30\textwidth}|p{0.26\textwidth}|p{0.30\textwidth}|}
\hline
\rowcolor{URJCRed}
\TVHead{Cambio} & \TVHead{Planteamiento inicial} & \TVHead{Enfoque final} \\
\hline
Papel de la Raspberry Pi & Posible elemento central de decisión y supervisión. & Nodo local de ingesta, almacenamiento, API y \textit{dashboard}. \\
\hline
Papel de ThingSpeak & Posible plataforma central de gestión. & \textit{Broker} de telemetría y consulta. \\
\hline
Papel de Arduino & Nodo de lectura de sensores. & Estación autónoma con lectura, decisión local y actuación segura. \\
\hline
Comunicación de campo & Dependencia potencial de red local o infraestructura específica. & Publicación LTE/HTTP desde la estación. \\
\hline
Actuación de riego & Posible control externo desde un nodo supervisor. & Control local temporizado desde el \textit{firmware}. \\
\hline
Validación & Comprobación directa del sistema completo. & Validación incremental por bloques para aislar fallos. \\
\hline
\end{tabular}
\caption{Cambios de enfoque relevantes durante el proyecto.}
\label{tab:cambios_enfoque}
\end{table}

\section{Lecciones aprendidas}

El seguimiento del proyecto ha permitido extraer varias lecciones técnicas:

\begin{itemize}
    \item En un sistema IoT físico, la arquitectura debe diseñarse desde el principio pensando en fallos de alimentación, sensores y comunicación.
    \item La actuación sobre una bomba debe apoyarse en lógica local, temporización y condiciones de seguridad verificables.
    \item Separar estación, telemetría, almacenamiento, API y \textit{dashboard} facilita la depuración del sistema.
    \item ThingSpeak resulta útil como \textit{broker} de telemetría y debe tratarse como una dependencia externa, manteniendo la lógica crítica de actuación en la estación.
    \item La conectividad LTE aporta valor en entornos donde se prioriza cobertura móvil, y exige cuidar alimentación, antena, SIM y configuración APN.
    \item La validación por bloques reduce errores acumulados y permite localizar mejor el origen de cada fallo.
\end{itemize}

\section{Conclusión}

El seguimiento del proyecto muestra que TierraViva ha evolucionado mediante decisiones técnicas sucesivas orientadas a mantener un equilibrio entre ambición, seguridad, reproducibilidad y alcance. La arquitectura resultante construye un prototipo coherente y funcional: estaciones autónomas que miden y actúan localmente, ThingSpeak como \textit{broker} de telemetría, Raspberry Pi como nodo de supervisión, SQLite como almacenamiento local y \textit{dashboard} web como interfaz de consulta.

Este anexo justifica qué se ha construido, por qué se ha construido de esa forma y qué alternativas quedaron fuera de la versión actual.